% Canonical TeX edition. Initial text inherited from 3b067cc2eba4ac2f2f91248448727639d8ac7c5a
% Source: research/case_study/phase238_topic_E_umbrella.md
% Subsequent mathematical edits require a new audit.
% Options for packages loaded elsewhere
\PassOptionsToPackage{unicode}{hyperref}
\PassOptionsToPackage{hyphens}{url}
\documentclass[
  11pt,
  a4paper,
]{article}
\usepackage{xcolor}
\usepackage[margin=1in]{geometry}
\usepackage{amsmath,amssymb}
\setcounter{secnumdepth}{-\maxdimen} % remove section numbering
\usepackage{iftex}
\ifPDFTeX
  \usepackage[T1]{fontenc}
  \usepackage[utf8]{inputenc}
  \usepackage{textcomp} % provide euro and other symbols
\else % if luatex or xetex
  \usepackage{unicode-math} % this also loads fontspec
  \defaultfontfeatures{Scale=MatchLowercase}
  \defaultfontfeatures[\rmfamily]{Ligatures=TeX,Scale=1}
\fi
\usepackage{lmodern}
\ifPDFTeX\else
  % xetex/luatex font selection
  \setmainfont[]{STIXGeneral}
  \setmonofont[]{Menlo}
\fi
% Use upquote if available, for straight quotes in verbatim environments
\IfFileExists{upquote.sty}{\usepackage{upquote}}{}
\IfFileExists{microtype.sty}{% use microtype if available
  \usepackage[]{microtype}
  \UseMicrotypeSet[protrusion]{basicmath} % disable protrusion for tt fonts
}{}
\makeatletter
\@ifundefined{KOMAClassName}{% if non-KOMA class
  \IfFileExists{parskip.sty}{%
    \usepackage{parskip}
  }{% else
    \setlength{\parindent}{0pt}
    \setlength{\parskip}{6pt plus 2pt minus 1pt}}
}{% if KOMA class
  \KOMAoptions{parskip=half}}
\makeatother
\usepackage{longtable,booktabs,array}
\newcounter{none} % for unnumbered tables
\usepackage{calc} % for calculating minipage widths
% Correct order of tables after \paragraph or \subparagraph
\usepackage{etoolbox}
\makeatletter
\patchcmd\longtable{\par}{\if@noskipsec\mbox{}\fi\par}{}{}
\makeatother
% Allow footnotes in longtable head/foot
\IfFileExists{footnotehyper.sty}{\usepackage{footnotehyper}}{\usepackage{footnote}}
\makesavenoteenv{longtable}
\setlength{\emergencystretch}{3em} % prevent overfull lines
\providecommand{\tightlist}{%
  \setlength{\itemsep}{0pt}\setlength{\parskip}{0pt}}
\usepackage{newunicodechar}
\newfontfamily\numeralfallback{STIX Two Text}
\newunicodechar{∎}{\ensuremath{\blacksquare}}
\usepackage{fvextra}
\DefineVerbatimEnvironment{Highlighting}{Verbatim}{breaklines,breakanywhere,breakanywheresymbolpre={},breakanywheresymbolpost={},breaksymbolleft={},breaksymbolright={},commandchars=\\\{\}}
\newunicodechar{⁰}{{\numeralfallback\char"2070\relax}}
\newunicodechar{₀}{{\numeralfallback\char"2080\relax}}
\newunicodechar{₁}{{\numeralfallback\char"2081\relax}}
\newunicodechar{₂}{{\numeralfallback\char"2082\relax}}
\newunicodechar{₃}{{\numeralfallback\char"2083\relax}}
\newunicodechar{₄}{{\numeralfallback\char"2084\relax}}

\usepackage{seqsplit}
\renewcommand{\texttt}[1]{{\ttfamily\seqsplit{#1}}}
\setlength{\emergencystretch}{6em}

\newunicodechar{⁻}{{\numeralfallback\char"207B\relax}}
\newunicodechar{⁸}{{\numeralfallback\char"2078\relax}}
\newunicodechar{ⁱ}{{\numeralfallback\char"2071\relax}}
\newunicodechar{⁴}{{\numeralfallback\char"2074\relax}}
\newunicodechar{⁷}{{\numeralfallback\char"2077\relax}}
\newunicodechar{⁶}{{\numeralfallback\char"2076\relax}}
\newunicodechar{⁵}{{\numeralfallback\char"2075\relax}}
\newunicodechar{₈}{{\numeralfallback\char"2088\relax}}
\newunicodechar{ᵢ}{\textsubscript{i}}
\newunicodechar{ⱼ}{\textsubscript{j}}
\newunicodechar{ₜ}{\textsubscript{t}}
\usepackage{bookmark}
\IfFileExists{xurl.sty}{\usepackage{xurl}}{} % add URL line breaks if available
\urlstyle{same}
% fallback for those not using the hyperref driver hyperxmp:
\makeatletter
\@ifundefined{xmpquote}{\newcommand{\xmpquote}[1]{#1}}{}
\makeatother
\hypersetup{
  hidelinks,
  pdfcreator={LaTeX via pandoc}}

\author{}
\date{}

\begin{document}

\section{Centralizers, facets and finite amalgams on the projective
line}\label{centralizers-facets-and-finite-amalgams-on-the-projective-line}

\subsection{Abstract}\label{abstract}

We give a common reduction for string C-groups of projective linear
degree two. A nonabelian dihedral subgroup defined over an
odd-characteristic subfield acquires no new centralizing projectivity
under field extension. Combined with an odd-order normalizer argument,
this excludes subfield facets without counting their mutual
intersections. The remaining local markings form six classes; their
fourteen compatible ordered pairs give ten universal presentations up to
reversal. We retain Dickson's subgroup types and the group--polytope
correspondence. The finite part is supplied by one exact computation.
Separate applications derive the complete projective special and general
linear principal classifications.

\subsection{1. The common theorem}\label{the-common-theorem}

Let \(q\) be a prime power. We write \(L(q)=\operatorname{PSL}_2(q)\)
and \(P(q)=\operatorname{PGL}_2(q)\). A string C-group is a group
\(G=\langle r_0,\ldots,r_{n-1}\rangle\) with distinct involutory
generators, with \(r_i r_j=r_j r_i\) when \(|i-j|>1\), and with the
intersection property \[
 G_I\cap G_J=G_{I\cap J},\qquad G_I=\langle r_i:i\in I\rangle.
 \tag{1}
\] The empty subgroup is trivial. In particular, no generator belongs to
the subgroup generated by the others. The type of a string C-group is
the sequence of orders of consecutive generator products.

\textbf{Theorem 1 (uniform reduction).} Suppose that a string C-group is
isomorphic to \(L(q)\) or \(P(q)\) and has rank at least four. Then
\(q\) is odd, the rank is four, and both end parabolics are proper and
isomorphic to \(S_4\) or \(A_5\). Their ordered markings are among the
six classes of Lemma 10. Up to reversal, the full marked group is a
quotient of one of the ten literal presentations of Lemma 11.

The assertion is a reduction, not yet a classification of all quotients
of those presentations. That last step belongs to the two application
papers. The common theorem includes small fields and arbitrary subfield
extension degrees. Its proof is completed in §6 after the general
centralizer argument and the finite table have been established.

The six local classes and the ten-row table are earlier results from
L083. The common proof below reworks their use in the two
classifications. The centralizer calculation and the normalizer assembly
replace the source's subfield-incidence counts. No claim of historical
priority for this proof method is needed for the applications.

\subsection{2. Retained inputs and elementary string
consequences}\label{retained-inputs-and-elementary-string-consequences}

We use two established background theorems.

\textbf{Input D (Dickson's classification, in its projective form).} For
a finite field \(K\) of odd characteristic, a proper subgroup of
\(L(K)\) or \(P(K)\) is one of the following types: cyclic; dihedral;
elementary abelian of characteristic order; a subgroup of an affine
group \(V\rtimes k^\times\); \(A_4\), \(S_4\), or \(A_5\); or a
conjugate of \(L(k)\) or \(P(k)\) for a subfield \(k\subseteq K\). Only
the types that actually embed are intended. The subfield alternatives
are conjugate to the usual matrix embeddings, not merely abstractly
isomorphic groups. We also retain the usual orders \[
 |L(q)|=\frac{q(q^2-1)}{(2,q-1)},\qquad |P(q)|=q(q^2-1),
 \tag{2}
\] the simplicity of \(L(q)\) for \(q\ge4\), and the trivial centralizer
of \(L(q)\) in \(P(q)\).

The required subgroup list is printed as Theorem 2 in L083 and as
Propositions 1--2 in L079, with attribution there to Dickson, Moore and
the later projective subgroup account. Their subgroup multiplicity
formulas are not inputs to our proof. In particular the affine
alternative includes dihedral groups of order \(2p\), where
\(p=\operatorname{char}K\); it must not be silently replaced by the
semisimple torus cases.

\textbf{Input P (the string C-group correspondence).} String C-markings,
up to group isomorphisms preserving the ordered marking, correspond to
regular abstract polytopes. The faces of rank \(i\) are the cosets of
the subgroup generated by all generators except \(r_i\). Reversing the
marking gives the dual polytope. This is the standard correspondence
used in both sources and in the existing proof library. A group
epimorphism preserving the ordered distinguished generators between two
string C-groups induces a covering of their regular polytopes. We retain
this covering part of the correspondence as well.

We first dispose of a possible degeneracy. If a consecutive product has
order two, all generators on one side of that place commute with all
generators on the other. By (1) the two generated subgroups meet
trivially. Since they generate \(G\), the multiplication map from their
direct product to \(G\) is an isomorphism. Thus a directly
indecomposable group has only types with consecutive entries at least
three.

All the groups \(L(q)\) and \(P(q)\) are directly indecomposable. For
\(q\ge4\), simplicity proves this for \(L(q)\). If \(P(q)=A\times B\),
its simple normal subgroup \(L(q)\) must lie in one factor: its
projections have commuting normal images, and a simple nonabelian group
cannot have two nontrivial such images. More explicitly, \(L(q)\cap A\)
and \(L(q)\cap B\) are normal, and if both were trivial, the commutators
of \(L(q)\) with both factors would be trivial. The other factor then
centralizes \(L(q)\), so is trivial. In even characteristic
\(P(q)=L(q)\). The remaining groups are \(S_3,A_4,S_4\), for which the
elementary normal-subgroup lists give the same conclusion. These lists
are \(1,A_3,S_3\); \(1,V_4,A_4\); and \(1,V_4,A_4,S_4\), respectively.

Every marking of rank at least four in these groups therefore has all
consecutive product orders at least three. Restricting a marking to an
interval of three indices gives an irreducible rank-three string
C-group.

\textbf{Lemma 2 (normalizing a facet).} Suppose
\(G=\langle H,s\rangle\), \(s^2=1\), and \(sHs^{-1}=H\). If \(t\in H\)
and \(ts\) has odd order, then \(s\in H\).

\textbf{Proof.} The normalizer assumption makes \(H\triangleleft G\).
The quotient is generated by \(sH\), whose order divides two. Since
\(tH=H\), the order of the same element \(tsH=sH\) divides the odd order
of \(ts\). Hence \(sH=H\). \(\square\)

This lemma is valid without a finiteness assumption on \(G\), provided
the displayed odd order exists. It is the common final step in both
subfield exclusions.

\subsection{3. Centralizers on the projective
line}\label{centralizers-on-the-projective-line}

We prove the matrix facts used in the reduction. All projectivities act
on the projective line; equality of matrices representing them is
understood up to nonzero scalar.

\textbf{Lemma 3 (characteristic two).} In \(P(K)\), for \(K\) of
characteristic two, the centralizer of every nonidentity involution is
abelian.

\textbf{Proof.} Over an algebraic closure an involution is unipotent,
since a semisimple element of order two would have eigenvalue ratio
\(-1=1\). Conjugate it to \(u:x\mapsto x+1\). It has the unique fixed
point infinity, so a commuting projectivity must fix infinity and has
the form \(x\mapsto ax+b\). Commutation gives \(a=1\). These
translations commute with one another. The centralizer over \(K\) is a
subgroup of this abelian group. \(\square\)

\textbf{Lemma 4 (commuting involutions in odd characteristic).} Let
\(K\) have odd characteristic. The centralizer of a four-group in
\(P(K)\) is that four-group. In particular \(P(K)\) contains no
elementary abelian subgroup of order eight.

\textbf{Proof.} We may extend scalars to an algebraic closure. Put the
first involution in the form \(z:x\mapsto-x\). A projectivity commuting
with \(z\) either preserves or interchanges its two fixed points, and
accordingly has the form \(x\mapsto ax\) or \(x\mapsto a/x\). The only
nontrivial involution of the first form is \(z\). An independent
commuting involution is therefore of the second form; rescaling the
coordinate makes it \(w:x\mapsto1/x\). The projectivities commuting with
both \(z\) and \(w\) are exactly \[
 x,\quad -x,\quad 1/x,\quad -1/x.
\] These are the given four-group. Since its elements were already
defined over \(K\), intersecting with \(P(K)\) changes nothing.
\(\square\)

Here is the field-extension statement responsible for the common
reduction.

\textbf{Theorem 5 (dihedral centralizer descent).} Let \(k\subseteq K\)
be fields of odd characteristic. Let \(a,b\in P(k)\) be distinct
involutions such that \(u=ab\) has finite order greater than two. Then
\[
 C_{P(K)}(\langle a,b\rangle)\ \subseteq\ P(k).
 \tag{3}
\] If \(u\) is unipotent, this centralizer is trivial. Otherwise it has
order two. In the semisimple case, if \(A\in\operatorname{GL}_2(k)\)
represents \(u\), its nontrivial element is represented by \[
 Z=A-\tfrac12\operatorname{tr}(A)I.
 \tag{4}
\]

\textbf{Proof.} First suppose \(u\) is semisimple. Over an algebraic
closure it has two distinct eigenlines, and we may write it as
\(x\mapsto\lambda x\), where \(\lambda\ne\pm1\). A projectivity
commuting with \(u\) must preserve its individual fixed points: an
interchange would conjugate \(u\) to \(u^{-1}\ne u\). Its form is
therefore \(x\mapsto cx\). Since \(aua^{-1}=u^{-1}\), the involution
\(a\) interchanges the eigenlines and has the form \(x\mapsto d/x\).
Commutation of \(x\mapsto cx\) with \(a\) gives \(c^2=1\). Thus the
centralizer over the algebraic closure is precisely \(\{1,z\}\), where
\(z:x\mapsto-x\).

Write the eigenvalues of \(A\) as \(\alpha,\beta\). The matrix in (4)
has eigenvalues \((\alpha-\beta)/2\) and \((\beta-\alpha)/2\). Both are
nonzero, and their ratio is \(-1\). Thus (4) represents \(z\), is
invertible, and has entries in \(k\). This proves both the exact order
and (3), including the nonsplit torus case. The expression is unaffected
projectively if \(A\) is rescaled.

Now suppose \(u\) is unipotent. Over the algebraic closure put
\(u:x\mapsto x+1\). As in Lemma 3 its centralizer consists of
translations \(x\mapsto x+c\). An involution conjugating \(u\) to
\(u^{-1}\) has form \(x\mapsto-x+d\). A translation commuting with it
satisfies \(2c=0\), hence \(c=0\). The centralizer is trivial. These two
cases exhaust nonidentity elements of \(P(k)\) in positive
characteristic. \(\square\)

The unipotent case explains directly why a dihedral group of order
\(2p\) cannot have an independent commuting involution in odd
characteristic. This is the case singled out in the corrective remark
following Proposition 2 of L079. No absence of dihedral groups of order
\(2p\) is assumed here.

\textbf{Corollary 6 (odd end labels).} In a rank-four string C-group
contained in \(P(K)\), with \(K\) of odd characteristic and all
consecutive labels at least three, the orders of \(r_0r_1\) and
\(r_2r_3\) are odd.

\textbf{Proof.} Put \(D=\langle r_0,r_1\rangle\) and let \(t=|r_0r_1|\).
The generator \(r_3\) centralizes \(D\) and lies outside it by (1).
Theorem 5 first excludes the unipotent case. If \(t\) were even, the
nontrivial involution \((r_0r_1)^{t/2}\in D\) would centralize \(D\).
The centralizer has exactly one nonidentity element, so it would equal
\(r_3\), a contradiction. Reverse the indices for the other end label.
\(\square\)

\subsection{4. The uniform rank and facet
reduction}\label{the-uniform-rank-and-facet-reduction}

\textbf{Proposition 7 (rank restriction).} A string C-group isomorphic
to \(L(q)\) or \(P(q)\) has rank at most four. If its rank is four,
\(q\) is odd.

\textbf{Proof.} In odd characteristic a marking of rank at least five
contains the independent commuting involutions \(r_0,r_2,r_4\). Their
generated subgroup is elementary abelian of order eight: (1) excludes
every dependence. This contradicts Lemma 4.

In characteristic two, a marking of rank at least four has
\(r_2,r_3\in C_G(r_0)\). Lemma 3 implies that they commute, contrary to
the consecutive label being at least three. This also excludes all
higher ranks. \(\square\)

\textbf{Lemma 8 (the nonexceptional rank-three groups).} An irreducible
rank-three string C-group cannot be cyclic, elementary abelian,
dihedral, \(A_4\), or a subgroup of the one-dimensional affine group in
odd characteristic.

\textbf{Proof.} Cyclic and elementary abelian groups are immediate. The
involutions of \(A_4\) generate its four-group and cannot generate
\(A_4\).

In an affine group in odd characteristic, an involution must have
multiplier \(-1\), so has form \(x\mapsto-x+v\). Two such involutions
commute only when their translation parameters agree. Thus there cannot
be two distinct commuting end generators.

For a dihedral group use its usual rotation/reflection description. A
central end generator would give an adjacent label two. Otherwise both
end generators are reflections. Two distinct commuting reflections have
product the central involution \(z\), so \(r_2=zr_0\). The middle
generator cannot be central, again by irreducibility, and hence is a
reflection. Put \(u=r_0r_1\). Then \[
 (r_1r_2)^2=(u^{-1}z)^2=u^{-2}.
\] The nonidentity rotation \(u^2\), nonidentity because \(|u|>2\),
belongs to both adjacent rank-two subgroups. It does not belong to the
reflection subgroup \(\langle r_1\rangle\). This violates (1).
\(\square\)

\textbf{Theorem 9 (exceptional facets).} Let \(G\cong L(q)\) or \(P(q)\)
have a rank-four string C-marking. Its facet and vertex-figure groups
are isomorphic to \(S_4\) or \(A_5\), and are proper subgroups.

\textbf{Proof.} Proposition 7 gives odd characteristic. Let
\(H=\langle r_0,r_1,r_2\rangle\). It is proper by (1) and irreducible by
§2. Input D and Lemma 8 leave only \(S_4,A_5\), and the conjugates of
standard subfield groups \(L(k)\) and \(P(k)\).

Suppose \(H\) is a subfield group, and conjugate the embedding so that
it is standard. The two involutions \(r_0,r_1\in H\subseteq P(k)\) have
product order greater than two. Theorem 5 places their centralizer in
\(P(k)\). Since \(r_3\) centralizes them, \(r_3\in P(k)\). If
\(H=P(k)\), this immediately contradicts (1). If \(H=L(k)\), the
determinant-square homomorphism shows that \(P(k)\) normalizes \(H\).
Hence \(r_3\) normalizes \(H\). Corollary 6 makes \(|r_2r_3|\) odd, so
Lemma 2, with \(t=r_2\), again gives \(r_3\in H\), a contradiction.
Reversing the marking proves the vertex-figure assertion. \(\square\)

No comparison of the extension degrees is needed. This proof includes
subfield embeddings of degree two, arbitrarily large extension degrees,
and the small subfield groups whenever the subgroup list permits them.
Exceptional groups abstractly isomorphic to \(P(3)\) or \(L(5)\) in a
\emph{different} characteristic are not treated as subfield embeddings;
they correctly remain among \(S_4,A_5\).

\subsection{5. The exceptional finite
calculation}\label{the-exceptional-finite-calculation}

For integers \(m,n,k\), write \([m,n]_k\) for the presentation \[
 \langle x,y,z\mid x^2=y^2=z^2=(xz)^2=(xy)^m=(yz)^n=(xyz)^k=1\rangle.
 \tag{5}
\] The corresponding map, when it is a string C-group, is traditionally
written \(\{m,n\}_k\). This is the Petrie relation convention of L083,
physical page 9.

\textbf{Finite Lemma 10.} Up to automorphisms preserving the ordered
marking, the irreducible rank-three string C-markings of \(S_4\) and
\(A_5\) have precisely the following triples \((m,n,k)\): \[
\begin{array}{c|ccc}
 S_4&(3,3,4)&(3,4,3)&(4,3,3)\\
 A_5&(3,5,5)&(5,3,5)&(5,5,3).
\end{array}
\tag{6}
\] In each case (5) is a presentation of the indicated group.

\textbf{Finite proof and audit procedure.} The program enumerates every
involution of \(S_4\) and \(A_5\), every ordered pair of distinct
commuting involutions for \(x,z\), and every middle involution \(y\). It
tests generation, both adjacent orders, and (1) on all eight subsets.
There are 72 surviving ordered triples in \(S_4\) and 360 in \(A_5\). It
partitions these by the full automorphism group, not merely inner
conjugation, and obtains the three ordered types in each row of (6). It
separately constructs (5) as a finitely presented group and checks order
24 or 60. A marking gives a surjection from (5) onto the corresponding
finite group; equality of orders makes it an isomorphism. These are
exhaustive finite calculations with GAP's exact subgroup and
finite-presentation machinery. \(\square\)

For two entries \(f=(m,n,k)\) and \(v=(n,s,l)\) from (6), define \[
\begin{split}
U(f,v)=\langle r_0,r_1,r_2,r_3\mid &r_i^2=1,\ (r_ir_j)^2=1\ (j-i>1),\\
 &(r_0r_1)^m=(r_1r_2)^n=(r_2r_3)^s=1,\\
 &(r_0r_1r_2)^k=(r_1r_2r_3)^l=1\rangle.
\end{split}\tag{7}
\] The marking of every rank-four polytope with these local maps is a
quotient of (7). Matching the common middle label in (6) gives fourteen
ordered pairs. Reversal sends \((f,v)\) to \((v^*,f^*)\), where
\((a,b,c)^*=(b,a,c)\). There are ten orbits, including six fixed pairs;
the following table contains one from each orbit.

{\def\LTcaptype{none} % do not increment counter
\begin{longtable}[]{@{}
  >{\raggedleft\arraybackslash}p{(\linewidth - 10\tabcolsep) * \real{0.2000}}
  >{\raggedright\arraybackslash}p{(\linewidth - 10\tabcolsep) * \real{0.1500}}
  >{\raggedright\arraybackslash}p{(\linewidth - 10\tabcolsep) * \real{0.1500}}
  >{\raggedleft\arraybackslash}p{(\linewidth - 10\tabcolsep) * \real{0.2000}}
  >{\raggedright\arraybackslash}p{(\linewidth - 10\tabcolsep) * \real{0.1500}}
  >{\raggedright\arraybackslash}p{(\linewidth - 10\tabcolsep) * \real{0.1500}}@{}}
\toprule\noalign{}
\begin{minipage}[b]{\linewidth}\raggedleft
Row
\end{minipage} & \begin{minipage}[b]{\linewidth}\raggedright
Facet \(f\)
\end{minipage} & \begin{minipage}[b]{\linewidth}\raggedright
Vertex figure \(v\)
\end{minipage} & \begin{minipage}[b]{\linewidth}\raggedleft
\(|U(f,v)|\)
\end{minipage} & \begin{minipage}[b]{\linewidth}\raggedright
Group
\end{minipage} & \begin{minipage}[b]{\linewidth}\raggedright
String C-marking
\end{minipage} \\
\midrule\noalign{}
\endhead
\bottomrule\noalign{}
\endlastfoot
1 & \((5,5,3)\) & \((5,5,3)\) & 1 & trivial & no \\
2 & \((5,5,3)\) & \((5,3,5)\) & 1 & trivial & no \\
3 & \((5,3,5)\) & \((3,5,5)\) & 3420 & \(L(19)\) & yes \\
4 & \((5,3,5)\) & \((3,4,3)\) & 60 & \(A_5\) & no \\
5 & \((5,3,5)\) & \((3,3,4)\) & 1 & trivial & no \\
6 & \((4,3,3)\) & \((3,4,3)\) & 96 & \(2^4:S_3\) & yes \\
7 & \((4,3,3)\) & \((3,3,4)\) & 24 & \(S_4\) & no \\
8 & \((3,5,5)\) & \((5,3,5)\) & 660 & \(L(11)\) & yes \\
9 & \((3,4,3)\) & \((4,3,3)\) & 1 & trivial & no \\
10 & \((3,3,4)\) & \((3,3,4)\) & 120 & \(S_5\) & yes \\
\end{longtable}
}

\textbf{Finite Lemma 11.} The table is correct. The groups in rows 3, 8
and 10 admit automorphisms reversing their four specified generators.
The face counts in these rows, from vertices through facets, are
respectively \[
 (57,171,171,57),\quad (11,55,55,11),\quad(5,10,10,5).
 \tag{8}
\] The graph of vertices and edges in row 8 is the complete graph
\(K_{11}\).

\textbf{Finite proof and audit procedure.} For each literal presentation
(7), the program computes the order of the finitely presented group,
constructs an isomorphic permutation group, and transports each
specified generator through that isomorphism. It tests (1) on all
sixteen subsets, including nonconsecutive ones. Thus a group
identification alone does not certify a polytope. The identity group is
handled separately, since GAP returns \texttt{fail} for its generic
finitely presented permutation-isomorphism request.

The three main rows are identified by an actual group isomorphism with
\(L(19),L(11),S_5\), respectively. Reversal is checked by constructing
the homomorphism on the four generators and checking bijectivity. The
program computes the four maximal-parabolic indices, giving (8). For row
8 it constructs the actual vertex and edge cosets and their incidences.
Every edge has two endpoints and the resulting 55 unordered pairs are
all distinct, proving the \(K_{11}\) assertion.

For the auxiliary row 6 it verifies an elementary abelian normal
subgroup of order 16, quotient \(S_3\), and the existence of a
complement. The indicated marking has face counts \((4,6,6,4)\). The
failures in rows 4 and 7 are checked by the intersection test;
independently, a stipulated facet or vertex-figure group already has the
order of the full group, so it cannot occur as the required proper
parabolic. The remaining failures are trivial groups. \(\square\)

This calculation recovers the actual ten-row table of L083, rather than
selecting only its two successful linear-group rows. The six local maps
and ten amalgams are prior results; the program supplies a new replay of
them.

\subsection{6. Proof of the reduction and
verification}\label{proof-of-the-reduction-and-verification}

\textbf{Proof of Theorem 1.} Section 2 excludes consecutive labels two,
including in the small groups. Proposition 7 then gives rank four in odd
characteristic. Theorem 9 identifies both proper end parabolics as
\(S_4\) or \(A_5\). Lemma 10 exhausts their ordered markings and
supplies an exact presentation for each local group. The consecutive
pairs share their middle product order. Substituting their two
presentations into the four-generator string relations gives the marked
epimorphism from (7). All compatible pairs, with their reversals, occur
in Lemma 11. This proves every assertion. \(\square\)

The common mathematical proof consists of this paper and the single
program \texttt{phase230\_finite\_check.g}. The finite-presentation
computations use GAP's exact algorithms to obtain the order of each
\emph{presented group}, and then an isomorphism with a permutation
group. Checking a chosen permutation image alone would give only a lower
bound and would not justify the table. Here the exact computation also
verifies all marked subgroup intersections, group identifications,
self-dualities and the listed coset incidences. The same program is used
by both applications; it is not rerun as an additional mathematical
hypothesis for either.

The inherited formal file
\texttt{phase230\_formal/DihedralFacetObstruction.lean} proves five
abstract steps: an involution with an odd vanishing power is trivial;
the corresponding homomorphism assertion; membership in a normal facet;
derivation of this conclusion from the normalizer condition and
generation; and contradiction with an independent generator. These are
parts of Lemma 2 and its use. They do not formalize the matrix
centralizers, Dickson's theorem, the finite presentations, or Theorem 1
as a whole. The finite program and those remaining written proofs retain
their stated verification status.

A portable replay verifies the six local classes and all ten universals.
Changing a Petrie relator is rejected by the universal-order check;
duplicating a distinguished generator is rejected by the full
intersection test. The formal replay reports only the usual Lean axioms.
The accompanying contract gives the exact files, versions and retained
source editions.

\subsection{7. Attribution and the application
interface}\label{attribution-and-the-application-interface}

Dickson's subgroup classification, the ordinary projective group
structure facts and the string C-group/regular-polytope correspondence
are retained inputs, not new results. The six exceptional rank-three
maps and the amalgam table are recovered from L083 with an explicit
replay. The classifications derived in the application papers belong to
Leemans and Schulte; the 11-cell and 57-cell themselves are the earlier
constructions of Grünbaum and Coxeter.

Relative to the existing proof bank, the work here is the explicit
projective-line centralizer descent and its common assembly with the
odd-order normalizer argument. An analogous centralizer strategy already
appears in the retained Suzuki discussion; it is not a substitute for
the matrix proof in §3. The present paper preserves the complete common
proof from the independently reviewed combined manuscript, while the two
applications carry their separate quotient and geometric obligations.

The applications use Theorem 1 and the exact marking data of Lemma 11.
Neither application is a premise of the other. Together they recover two
principal classification contracts; they do not reprove every secondary
statement in the historical sources. In particular the almost-simple
seventeen-case census and the separate conjecture in L079 are not
consequences asserted here. The topic's cost includes this entire common
paper and its required finite program once, followed by each complete
application. A common lemma used in another topic is charged again as
part of that other topic's full chosen umbrella.

\subsection{References}\label{references}

{[}L083{]} D. Leemans and E. Schulte, \emph{Groups of type \(L_2(q)\)
acting on polytopes}, retained arXiv:math/0606660v1 author edition, 14
physical pages; \emph{Advances in Geometry} 7 (2007), 529--539. Theorem
2 prints the retained subgroup classification; §5 prints the exceptional
amalgam table.

{[}L079{]} D. Leemans and E. Schulte, \emph{Polytopes with groups of
type \(PGL_2(q)\)}, retained arXiv:0909.1991v1, 9 physical pages.
Propositions 1--2 print the projective subgroup inputs; Theorem 10 is
the principal projective general linear classification.

{[}MS{]} P. McMullen and E. Schulte, \emph{Abstract Regular Polytopes},
Cambridge University Press, 2002. The marked-group correspondence, face
cosets, duality and marked epimorphism/covering theorem are retained.

\end{document}
